\section{Introduction}

Autonomous UAVs must often route while obstacles, regulatory restrictions, and
energy conditions change during flight. A route that is valid at takeoff can
become inefficient or infeasible after deployment, making dynamic replanning a
central navigation problem \cite{ghambari2024survey}.

Classical graph search provides strong reference methods when the current map
is known. Dijkstra is optimal but expands uniformly; A* preserves optimality
while using a heuristic to focus search; D* Lite reuses prior search state when
edge availability changes \cite{hart1968astar,koenig2005dstarlite}. These three
algorithms represent meaningfully different classical regimes: naive full
replanning, heuristic full replanning, and incremental replanning.

Deep reinforcement learning (RL) offers a different deployment model. After
training, a policy selects actions from observations without explicitly
solving a graph-search problem at every decision. This may be attractive under
rapid changes or large state spaces, but it removes shortest-path guarantees
and introduces sensitivity to training seed, representation, reward design,
and distribution shift \cite{vanhasselt2016double,andrychowicz2017her,
henderson2018deep}.

A credible comparison must therefore address more than one policy and one
classical baseline. It must pair methods on identical scenarios, retain
training and timing distributions, test unseen environments, separate
in-distribution from out-of-distribution performance, and measure adaptation
after individual change events. This paper implements that evidence standard
for a controlled grid-based UAV-routing study.

The study asks four questions:
\begin{enumerate}
    \item How do RL, naive search, heuristic search, and incremental search
    compare in reliability, path quality, and route-level decision cost?
    \item How much does RL performance vary across independent training seeds?
    \item How do the methods respond to layout shift, denser obstacles, new
    dynamics, realistic cost factors, and increasing grid size?
    \item Which RL components improve performance under controlled
    one-factor interventions?
\end{enumerate}

The principal contributions are:
\begin{itemize}
    \item a persisted 310-scenario benchmark that pairs methods on identical
    graphs, changes, movement costs, information timing, and compute boundaries;
    \item a 35-run, 17.5-million-step training study covering five seeds and
    controlled interventions on Double DQN, HER, reward shaping, curriculum,
    observation scope, and dynamic-only training;
    \item seed-aware generalization, scaling, realism, timing, and event-level
    adaptability analyses with paired tests and multiplicity correction; and
    \item an auditable pipeline linking source provenance, checkpoints, raw
    repetitions, tables, vector figures, and manuscript claims.
\end{itemize}
